\begin{center}\bfseries 10.\end{center}
\noindent\hspace*{\parindent}Zu den früheren Betrachtungen zurückkehrend begegnen wir der Frage nach dem Verhältnis der Anzahl der Klassen in irgend einer Ordnung \(O\) zur Anzahl der eigentlich primitiven Klassen. Es findet sich gleich der Anzahl der letzteren Formen, welche mit einer besonderen Form jener Ordnung zusammengesetzt diese Form reproduzieren. Heisst diese \((A,B,C)\), so sind das diejenigen eigentlich primitiven Formen, durch welche \(A^2\) darstellbar ist. \textsc{Gauss} hat ihre Anzahl berechnet für negative Determinanten, für positive aber nur zwischen ihr und der Fundamentalauflösung der \textsc{Pell}schen Gleichung einen Zusammenhang erkannt, den erst nach ihm \textsc{Dirichlet} klargelegt und bestimmt hat\footnote{\textit{Recherches sur diverses applications de l’Analyse infinitésimale à la théorie des nombres} § 8, \textsc{Crelles} Journal für Mathematik \textbf{21}, 1840, S. 10, Werke I, S. 470.}.

Ferner heben wir den Satz hervor, dass in jedem Geschlechte derselben Ordnung gleichviel Klassen befindlich sind, beschränken fortan die Betrachtung auf die eigentlich primitive Ordnung und müssen nun besonders der sogenannten Anzepsklassen gedenken, die dadurch charakterisiert sind, dass sie durch Duplikation d. i. durch Zusammensetzung mit sich selbst die Hauptklasse hervorbringen. Von grösster Wichtigkeit ist die \textsc{Gauss}sche Bestimmung ihrer Anzahl \(\alpha\), welche gleich der halben Anzahl aller angebbaren Gesamtcharaktere: \(\alpha=\frac12\chi\) gefunden wird. Andererseits ist die Anzahl aller eigentlich primitiven Klassen \(h=g\mathbin{.}\gamma\), wo \(\gamma\) die Anzahl ihrer Geschlechter, während sich zeigen lässt, dass auch \(h=\alpha\mathbin{.}\delta\) gesetzt werden kann, wenn \(\delta\) die Anzahl derjenigen unter jenen Klassen bedeutet, welche durch Duplikation entstehen können. Da alle diese aber zum Hauptgeschlechte gehören müssen, sodass \(\delta\EqualAboveLess g\) sein muss, so ergibt sich \(\gamma\EqualAboveLess\alpha\) d. h. der Satz, dass die Anzahl der Geschlechter höchstens halb so gross, wie die aller angebbaren Gesamtcharaktere, dass also für die Hälfte der letztern gewiss kein entsprechendes Geschlecht vorhanden ist.

Nur beiläufig und doch als ein Glanzpunkt in der Reihe dieser Betrach-\newpage
\SourceLeaf{32}\ArithmeticHead{25}{R}{TERNÄRE QUADRATISCHE FORMEN.\quad DER VI. ABSCHNITT DER D. A.}\HistoricalFoot{X 2.}{4}
\noindent tungen knüpft sich an das letztere Ergebnis, indem es auf den Fall angewandt wird, wo die Determinante eine positive oder negative Primzahl bezw. das Produkt von zwei solchen ist, der zweite, am 27. Juni 1796 (Nr. 16 T.) gefundene \textsc{Gauss}sche Beweis des Reziprozitätsgesetzes und seiner beiden Ergänzungssätze an, ein Beweis, der, in seiner Grundlage scheinbar so gänzlich fremd der Frage, auf welche er zielt, eben wegen der Verbindung ganz verschiedener Gedankenreihen, die er knüpft, als der tiefste all’ seiner Beweise des Gesetzes bezeichnet werden darf.
